\section{Results}
\label{sec:results}

\subsection{Unified benchmark and cross-backbone comparison}

We report PR-AUC and ROC-AUC for binary tasks and Cohen's $\kappa$ and weighted
F1 for multiclass tasks. Table~\ref{tab:unified-results} gives a unified
comparison between the three visual backbones and the available EEG
foundation-model references (BIOT, LaBraM-Base, CBraMod, and REVE). Balanced
accuracy, population standard deviations, and detailed supervised baselines
are provided in the appendix.

\begin{table*}[t]
\centering
\caption{Unified comparison across 12 downstream datasets. Binary tasks report
PR-AUC / ROC-AUC; multiclass tasks report Cohen's $\kappa$ / weighted F1.
Bold and underlined values indicate the best and second-best results. A
superscript $^*$ marks datasets where at least one visual model exceeds all
available EEG foundation-model references; missing entries are excluded.}
\label{tab:unified-results}
\tiny
\renewcommand{\arraystretch}{1.12}
\setlength{\tabcolsep}{1.5pt}
\begin{tabularx}{\textwidth}{@{}p{1.45cm}*{7}{Y}@{}}
\toprule
Dataset & \multicolumn{4}{c}{\textit{EEG foundation models}} & \multicolumn{3}{c}{\textit{Vision models}} \\
\cmidrule(lr){2-5}\cmidrule(l){6-8}
 & BIOT & LaBraM & \cbramod{} & \reve{} & EffNet-B0 & ConvNeXt-T & ViT-S \\
\midrule
CHB-MIT$^*$ & .3277/.8761 & .3287/.8679 & .3689/.8892 & -- & .4092/.8974 & \underline{.4734}/\underline{.9067} & \textbf{.4759}/\textbf{.9116} \\
TUAB & .8792/.8815 & .8965/.9022 & \underline{.9221}/\underline{.9156} & \textbf{.9281}/\textbf{.9245} & .9178/.9126 & .9057/.9011 & .9091/.9042 \\
SHU-MI$^*$ & .6770/.6609 & .6761/.6604 & \textbf{.7139}/\underline{.6988} & -- & \underline{.6982}/.6962 & .6910/\textbf{.7006} & .6891/.6908 \\
Mumtaz2016 & .9736/.9758 & .9798/.9782 & \underline{.9923}/\underline{.9921} & \textbf{.9961}/\textbf{.9957} & .9864/.9846 & .9842/.9825 & .9797/.9767 \\
MentalArithmetic$^*$ & .6004/.7536 & .5999/.7721 & .6267/.7905 & .7470/.8450 & \textbf{.7903}/\textbf{.8779} & \underline{.7598}/\underline{.8752} & .7249/.8498 \\
\midrule
TUEV$^*$ & .5273/.7492 & .6637/.8312 & .6744/.8331 & .6783/.8451 & .6901/.8380 & \underline{.7095}/\underline{.8488} & \textbf{.7428}/\textbf{.8655} \\
ISRUC$^*$ & .7192/.7790 & .7231/.7810 & .7442/.8011 & .7500/.8005 & .7705/.8197 & \textbf{.7906}/\textbf{.8372} & \underline{.7875}/\underline{.8349} \\
HMC$^*$ & .6295/.7091 & .6812/.7554 & -- & .6982/.7638 & \underline{.7070}/\underline{.7700} & \textbf{.7111}/\textbf{.7744} & .6987/.7634 \\
FACED$^*$ & .4476/.5136 & .4698/.5288 & .5041/.5618 & \underline{.5080}/\underline{.5659} & .4953/.5567 & \textbf{.5891}/\textbf{.6374} & .3744/.4445 \\
SEED-V & .2261/.3856 & .2386/.3974 & \textbf{.2569}/\textbf{.4101} & -- & .2056/.3678 & .2396/.3976 & \underline{.2442}/\underline{.4016} \\
PhysioNet-MI$^*$ & .4875/.6158 & .4912/.6177 & .5222/.6427 & .5306/.6484 & \underline{.5322}/\underline{.6485} & \textbf{.5479}/\textbf{.6618} & .5169/.6386 \\
BCIC2020-3$^*$ & .3650/.4917 & .3800/.5054 & .4216/.5383 & .4543/.5633 & \underline{.5673}/\underline{.6539} & .5560/.6447 & \textbf{.5773}/\textbf{.6618} \\
\bottomrule
\end{tabularx}
\tablenote{Published baselines are transcribed from \cbramod{} and \reve{} under the shared benchmark splits; local vision rows are five-seed means.}
\end{table*}

The unified table shows a heterogeneous but systematic pattern. At least one
visual backbone exceeds every displayed EEG foundation-model reference on nine
datasets, including CHB-MIT, TUEV, ISRUC, FACED, PhysioNet-MI, and BCIC2020-3.
The result is not tied to one encoder: ViT-Small is strongest on CHB-MIT, TUEV,
and BCIC2020-3; ConvNeXt-Tiny on ISRUC, HMC, FACED, and PhysioNet-MI; and
EfficientNet-B0 on MentalArithmetic. Thus the interface transfers across both
CNN and Transformer backbones, although backbone choice remains task-dependent.

The comparison also defines clear limits. REVE remains strongest on TUAB and
Mumtaz2016, while CBraMod leads on SEED-V and SHU-MI for at least one metric;
visual transfer is therefore competitive, not a universal replacement for EEG
pretraining. PhysioNet-MI further shows that competitive transfer can occur at
$P=1$, without folding. The TUEV release difference leaves the held-out test
set unchanged. Detailed supervised controls are in Appendix~\ref{app:detailed-results};
published foundation-model rows are contextual references rather than local
reruns.

\subsection{Matched initialization ablation}
\label{sec:pretraining-ablation}

The random-initialization controls test whether the visual weights themselves
transfer to EEG. Each comparison keeps the architecture, folding, optimizer,
epoch budget, checkpoint selection, and five-seed test protocol fixed; only
the backbone initialization changes. Table~\ref{tab:vision-pretraining-ablation}
reports the four representative datasets TUEV, PhysioNet-MI, ISRUC, and SHU-MI;
the full 12-dataset by 3-backbone sweep is given in Appendix~\ref{app:complete-ablations}.
Balanced accuracy is reported as the unified comparison metric.

\begin{table*}[t]
\centering
\caption{Matched visual-pretraining ablation on representative datasets.
Values are balanced-accuracy mean$\pm$population standard deviation; $\Delta$
is pretrained minus random, reported in percentage points.}
\label{tab:vision-pretraining-ablation}
\small
\setlength{\tabcolsep}{7pt}
\renewcommand{\arraystretch}{1.12}
\begin{tabular}{@{}llrrr@{}}
\toprule
Backbone & Dataset & Pretrained & Random & $\Delta$ (pp) \\
\midrule
EfficientNet-B0 & TUEV & .6423$\pm$.0164 & .5328$\pm$.0357 & +10.95 \\
 & PhysioNet-MI & .6491$\pm$.0125 & .5931$\pm$.0138 & +5.60 \\
 & ISRUC & .8025$\pm$.0027 & .7989$\pm$.0030 & +0.37 \\
 & SHU-MI & .6273$\pm$.0085 & .5005$\pm$.0004 & +12.68 \\
\midrule
ConvNeXt-Tiny & TUEV & .6763$\pm$.0122 & .4302$\pm$.0300 & +24.60 \\
 & PhysioNet-MI & .6609$\pm$.0079 & .4523$\pm$.0089 & +20.86 \\
 & ISRUC & .8142$\pm$.0009 & .6807$\pm$.0207 & +13.34 \\
 & SHU-MI & .6358$\pm$.0075 & .5775$\pm$.0093 & +5.83 \\
\midrule
ViT-Small & TUEV & .7261$\pm$.0135 & .4902$\pm$.0154 & +23.59 \\
 & PhysioNet-MI & .6377$\pm$.0042 & .5774$\pm$.0060 & +6.03 \\
 & ISRUC & .8130$\pm$.0027 & .7675$\pm$.0036 & +4.55 \\
 & SHU-MI & .6302$\pm$.0331 & .5864$\pm$.0115 & +4.38 \\
\bottomrule
\end{tabular}
\end{table*}

Across the complete 36 matched backbone--dataset comparisons, pretraining
improves balanced accuracy in 34 cases; complete values are reported in
Appendix~\ref{app:complete-ablations}. The representative rows show large
gains on TUEV and PhysioNet-MI, a small gain on ISRUC, and a
backbone-dependent result on SHU-MI. The magnitude is therefore task- and
backbone-dependent. The smallest gains occur on
large or easy datasets, whereas the
largest gains occur on FACED, BCIC2020-3, and MentalArithmetic. Thus generic visual
representations provide a reusable initialization for EEG fine-tuning after
geometry alignment, and the pretraining benefit grows with model capacity.
This is evidence for transfer from visual pretraining to EEG, not evidence
that EEG pretraining is unnecessary in frozen, few-shot, or cross-montage
settings.

\subsection{Phase-interleaved folding versus contiguous reshape}
\label{sec:fold-reshape}

To isolate temporal adjacency from simple height expansion, we compare
phase-interleaved folding with a contiguous-chunk reshape. Both produce the
same $CP\times(T/P)$ tensor and retain the same samples; only their time-to-row
mapping differs. Across the three-backbone controls, phase folding is favored
in 26 of 27 primary-metric comparisons: it is favored in all 18 CNN comparisons and in 8 of
9 ViT comparisons, with one small CHB-MIT reversal. The clearest gains occur
on FACED, MentalArithmetic, SHU-MI, and HMC. Table~\ref{tab:fold-representative}
shows representative primary-metric values for the two CNN backbones; the
complete backbone-specific results are reported in Appendix~\ref{app:fold-reshape}.
Thus the comparison tests the
ordering of samples presented to the backbone, not merely the increase in
input height.

\begin{table}[tbp]
\centering\small
\caption{Representative phase-interleaved folding gains. Values are primary-metric means
$\pm$ population standard deviation to four decimal places; $\Delta$ is phase fold minus contiguous reshape in
percentage points.}
\label{tab:fold-representative}
\begin{tabular}{@{}llrrr@{}}
\toprule
Backbone & Dataset / metric & Temporal fold & Contiguous reshape & $\Delta$ (pp) \\
\midrule
EffNet-B0 & FACED / $\kappa$ & .4953$\pm$.0197 & .3198$\pm$.0649 & +17.55 \\
 & MentalArithmetic / PR-AUC & .7903$\pm$.0536 & .5849$\pm$.0734 & +20.54 \\
 & HMC / $\kappa$ & .7070$\pm$.0023 & .6379$\pm$.0046 & +6.91 \\
 & SHU-MI / PR-AUC & .6982$\pm$.0178 & .7009$\pm$.0138 & $-0.27$ \\
\midrule
ConvNeXt-T & FACED / $\kappa$ & .5890$\pm$.0032 & .3321$\pm$.1694 & +25.69 \\
 & MentalArithmetic / PR-AUC & .7598$\pm$.0436 & .6198$\pm$.1237 & +14.01 \\
 & HMC / $\kappa$ & .7111$\pm$.0063 & .6700$\pm$.0050 & +4.11 \\
 & SHU-MI / PR-AUC & .6910$\pm$.0129 & .6822$\pm$.0333 & +0.88 \\
\bottomrule
\end{tabular}
\end{table}

\subsection{Fold-factor sensitivity}
\label{sec:fold-factor-sensitivity}

The fold factor is itself tunable. Table~\ref{tab:p-representative} summarizes
three representative backbone--dataset sweeps; complete results are given in
Appendix~\ref{app:p-ablation}. Across the evaluated sweeps, the
geometry-based default lies near the best-performing region despite using no
task-specific test-set tuning. The preferred value nevertheless varies with
task, backbone, and metric: larger factors are favored for CHB-MIT, whereas
TUEV and MentalArithmetic favor different settings across backbones. This
suggests that validation-based per-dataset selection of $P$ could yield
additional gains. We therefore use the geometry rule as a reproducible default
rather than a universal optimum, and select any tuned value using validation
data only.

\begin{table}[tbp]
\centering\small
\caption{Representative fold-factor sensitivity. We show one representative sweep for each backbone; entries are primary-metric
means $\pm$ population standard deviation to four decimal places; bold marks
the best tested factor.}
\label{tab:p-representative}
\begin{tabular}{@{}lrrrr@{}}
\toprule
Dataset / backbone & $P=1$ & $P=2$ & $P=4$ & $P=8$ \\
\midrule
CHB-MIT / EfficientNet-B0 (PR-AUC) & .3686$\pm$.0958 & .3676$\pm$.0498 & \textbf{.4092$\pm$.0536} & .3879$\pm$.0493 \\
TUEV / ConvNeXt-Tiny ($\kappa$) & .5381$\pm$.1921 & .7042$\pm$.0414 & .7095$\pm$.0305 & \textbf{.7274$\pm$.0124} \\
MentalArithmetic / ViT-Small (PR-AUC) & .7814$\pm$.0259 & \textbf{.7975$\pm$.0383} & .7249$\pm$.0625 & .5565$\pm$.0394 \\
\bottomrule
\end{tabular}
\end{table}

\subsection{Model complexity}
\label{sec:model-complexity}

Table~\ref{tab:model-scale} compares encoder parameters and forward compute
on a shared $32\times1000$ EEG input with $P=4$ for visual models.
EfficientNet-B0 uses 4.01M parameters and 0.247 GFLOPs, compared with 4.92M
and 0.79 GFLOPs for \cbramod{} and
69.19M and 11.07 GFLOPs for \reve{}-Base. ConvNeXt-Tiny and ViT-Small require
2.903 and 3.000 GFLOPs, respectively. Folding introduces no learned parameters
or arithmetic operations, but requires $O(CT)$ memory movement. These counts
characterize model complexity rather than measured runtime or memory use.
